% Overview of the study design: pipeline, validate-and-retry loop, and the
% two-stage classifier. Requires: tikz with \usetikzlibrary{positioning,arrows.meta}, adjustbox
\begin{figure}[!htbp]
\centering
\begin{adjustbox}{max width=\linewidth}
\begin{tikzpicture}[
  font=\footnotesize,
  >={Stealth[length=2mm]},
  node distance=4mm and 6mm,
  box/.style={draw, rounded corners=2pt, align=center, inner sep=3pt,
              minimum height=7mm},
  node/.style={box, fill=black!6, minimum width=17mm},
  stp/.style={box, minimum width=24mm},
  cat/.style={align=left, inner sep=2pt},
  stage/.style={draw, rounded corners=2pt, inner sep=4pt, align=left},
]
% --- pipeline row
\node[box, minimum width=14mm] (q) {Question};
\node[node, right=of q] (n1) {\textbf{identify}};
\node[node, right=of n1] (n2) {\textbf{gather}};
\node[node, right=of n2] (n3) {\textbf{answer}};
\node[box, right=of n3, minimum width=14mm] (out) {Result};
\draw[->] (q) -- (n1);
\draw[->] (n1) -- (n2);
\draw[->] (n2) -- (n3);
\draw[->] (n3) -- (out);

% --- one attempt at a node
\node[stp, below=9mm of n1] (prompt) {Prompt with\\JSON Schema};
\node[stp, right=of prompt] (llm) {Model\\completion $y$};
\node[stp, right=of llm] (val) {Pydantic\\validation};
\draw[->] (prompt) -- (llm);
\draw[->] (llm) -- (val);
\draw[->, dashed] (n2.south) -- node[right, pos=0.45, font=\scriptsize]{each node} (llm.north);
\draw[->] (val.east) -- ++(8mm,0) node[right, font=\scriptsize, align=left]{valid:\\next node};
\draw[->] (val.south) -- ++(0,-4mm) -| node[pos=0.25, below, font=\scriptsize]
  {invalid: retry with validation error appended (up to 3 further attempts)} (prompt.south);

% --- classifier
\node[stage, below=15mm of prompt.south west, anchor=north west, text width=55mm] (s1)
  {\raggedright\textbf{Stage 1:} predicates on raw text $y$, in order\\[2pt]
   \cat{markdown\_wrap}: starts with a \texttt{\textasciigrave\textasciigrave\textasciigrave} fence\\
   \cat{prose\_preamble}: prose before the JSON\\
   \cat{schema\_mirroring}: returns schema keywords\\[2pt]
   \emph{first match wins; precedence over Stage 2}};
\node[stage, right=6mm of s1.north east, anchor=north west, text width=55mm] (s2)
  {\raggedright\textbf{Stage 2:} if no predicate fires,\\map the Pydantic error type\\[2pt]
   \cat{regex\_pattern}, \cat{length\_bound}, \cat{enum\_violation},
   \cat{missing\_field}, \cat{range\_bound}, \cat{parse\_error},
   \cat{cross\_field\_ref}, catch-all\\[1pt]
   several errors: plurality category};
\draw[->] (s1.east |- s2.north west) ++(0,-3mm) -- ++(6mm,0);
\node[font=\scriptsize, anchor=south west] at (s1.north west) {Classifier: each failed attempt receives one primary category};
\end{tikzpicture}
\end{adjustbox}
\caption{Study design. Each question passes through three nodes, each with
its own Pydantic schema stated in the prompt; a failed validation is
retried with the error appended. Every failed attempt is assigned one
primary category, first by behavioural predicates on the raw completion
and otherwise by the Pydantic error type (\cref{sec:setup:taxonomy}).}
\label{fig:overview}
\end{figure}
